\section{Интеграция и работа модуля в системе}
\label{sec:work}

Агент интегрирован в общую программу обработки сообщений NIKA. Python-модуль \texttt{MessageProcessingModule} регистрирует \texttt{MusicEntityEnrichmentAgent} наравне с агентом прогноза погоды:

\begin{lstlisting}[caption={Регистрация агента в Python-модуле},label={lst:module},basicstyle=\footnotesize\ttfamily]
class MessageProcessingModule(ScModule):
    def __init__(self):
        super().__init__(WeatherAgent(), MusicEntityEnrichmentAgent())
\end{lstlisting}

Класс действия \texttt{action\_enrich\_music\_entity} описан в БЗ как атомарное действие класса \texttt{information\_action} с русским и английским идентификаторами (см. раздел~\ref{sec:agent}).

Вызов обогащения встроен в агента стандартного ответа (\texttt{StandardMessageReplyAgent.cpp}): если сообщение~\mdash  о сущности и у сущности нет \texttt{nrel\_note}, создаётся действие и ожидается его завершение:

\begin{lstlisting}[caption={Фрагмент StandardMessageReplyAgent: инициирование обогащения},label={lst:cpp},basicstyle=\footnotesize\ttfamily]
void StandardMessageReplyAgent::enrichMusicEntityIfNeeded(
    const ScAddr & messageNode)
{
  // is the message about an entity (concept_message_about_entity),
  // does it carry an entity via rrel_entity, and is nrel_note missing?
  if (!isAboutEntityWithMissingNote(messageNode))
    return;
  ScAddr enrichAction = AgentUtils::initAgent(
      &m_memoryCtx, MessageKeynodes::action_enrich_music_entity,
      {messageNode});
  ActionUtils::waitAction(&m_memoryCtx, enrichAction, 15000);
}
\end{lstlisting}

Ключевой фрагмент Python-агента~\mdash  запрос к MediaWiki API и запись результата в БЗ:

\begin{lstlisting}[caption={Фрагмент MusicEntityEnrichmentAgent: поиск и запись},label={lst:py},basicstyle=\footnotesize\ttfamily]
search_params = {"action": "query", "list": "search",
                 "srsearch": entity_name, "format": "json"}
resp = requests.get(WIKI_API, params=search_params, timeout=15)
page_title = resp.json()["query"]["search"][0]["title"]
extract_params = {"action": "query", "prop": "extracts",
                  "exintro": True, "explaintext": True,
                  "titles": page_title, "format": "json"}
summary = requests.get(WIKI_API, params=extract_params,
                       timeout=15).json() ... ["extract"]

note_link = create_link(summary, sc_types.LINK_CONST)
note_edge = create_edge(sc_types.EDGE_D_COMMON_CONST,
                        entity_addr, note_link)
create_edge(sc_types.EDGE_ACCESS_CONST_POS_PERM,
            nrel_note, note_edge)
create_edge(sc_types.EDGE_ACCESS_CONST_POS_PERM,
            lang_ru, note_link)
\end{lstlisting}
% The two create_edge calls above attach the note text to the entity
% under nrel_note and mark the link as Russian (lang_ru).

Классификация сообщения выполняется модулем witai-shim: намерение \texttt{about\_entity} определяется по шаблонам <<что такое>>, <<кто такой>>, <<расскажи о ...>> из файла \texttt{classification/intents.json}, сущность распознаётся по словарю \texttt{concept\_wit\_entity} предметной области.

На основе погруженных знаний система выводит ответ: логическое правило \texttt{lr\_message\_about\_entity} находит сущность с \texttt{nrel\_note}, а шаблон фразы \texttt{concept\_phrase\_about\_entity} подставляет имя и заметку:

\begin{lstlisting}[caption={Шаблон фразы ответа},label={lst:phrase},basicstyle=\footnotesize\ttfamily]
${_entity_note_name} - ${_note}.
\end{lstlisting}

Полный поток обработки сообщения показан на рисунке~\ref{fig:flow}.

\begin{figure}[H]
	\centering
	\begin{tikzpicture}[
		font=\footnotesize,
		node distance=0.55cm,
		box/.style={rectangle, rounded corners, minimum width=6.4cm, minimum height=0.8cm, align=center, draw=black, fill=white},
		agent/.style={rectangle, rounded corners, minimum width=6.4cm, minimum height=0.8cm, align=center, draw=black, fill=gray!12},
		ext/.style={rectangle, rounded corners, minimum width=4.0cm, minimum height=0.8cm, align=center, draw=black, fill=blue!8},
		arrow/.style={-{Stealth[length=2.1mm]}, semithick}
	]
		\node[box] (user) {Пользователь};
		\node[agent, below=of user] (cls) {Классификатор сообщений (witai-shim):\\намерение \texttt{about\_entity} + сущность};
		\node[agent, below=of cls] (std) {StandardMessageReplyAgent (C++):\\класс сообщения, \texttt{rrel\_entity}, проверка \texttt{nrel\_note}};
		\node[agent, below=of std] (music) {MusicEntityEnrichmentAgent (Python):\\найти имя, запросить статью, записать заметку};
		\node[ext, right=1.3cm of music] (wiki) {MediaWiki API\\(Википедия)};
		\node[agent, below=of music] (kb) {БЗ: сущность\\+ \texttt{nrel\_note}};
		\node[agent, below=of kb] (rule) {Правило и фраза \texttt{lr\_message\_about\_entity}};
		\node[box, below=of rule] (reply) {Ответ в интерфейсе NIKA};

		\draw[arrow] (user) -- (cls);
		\draw[arrow] (cls) -- (std);
		\draw[arrow] (std) -- node[anchor=east] {нет \texttt{nrel\_note}} (music);
		\draw[arrow] (music) -- (wiki);
		\draw[arrow] (wiki) -- (music);
		\draw[arrow] (music) -- (kb);
		\draw[arrow] (kb) -- (rule);
		\draw[arrow] (rule) -- (reply);
	\end{tikzpicture}
	\caption{Иллюстрация работы модуля адаптации знаний в системе NIKA}
	\label{fig:flow}
\end{figure}
